\documentclass[11pt,a4paper]{article}
\usepackage[T1]{fontenc}
\usepackage[utf8]{inputenc}
\usepackage[french]{babel}
\usepackage[margin=2cm]{geometry}
\usepackage{array}
\usepackage{tabularx}
\usepackage{booktabs}
\usepackage{colortbl}
\usepackage{xcolor}
\usepackage{ragged2e}
\usepackage{enumitem}
\usepackage{fancyhdr}

\definecolor{entete}{RGB}{31,73,125}
\definecolor{ligne}{RGB}{234,240,248}

\newcolumntype{L}[1]{>{\RaggedRight\arraybackslash}p{#1}}
\newcommand{\dimtitle}[1]{\vspace{4pt}\noindent{\color{entete}\large\bfseries #1}\par\vspace{2pt}}

\pagestyle{fancy}\fancyhf{}
\lhead{\small Construction des indicateurs N-MODA}
\rhead{\small EHCVM 2018/19 -- 2021/22}
\cfoot{\thepage}
\renewcommand{\headrulewidth}{0.4pt}

\setlength{\parindent}{0pt}
\renewcommand{\arraystretch}{1.25}

\begin{document}

\begin{center}
  {\LARGE\bfseries\color{entete} Construction des indicateurs de pauvreté\\[2pt]
   multidimensionnelle de l'enfant (N-MODA)}\\[6pt]
  {\normalsize Sié Rachid \textsc{Traoré} --- ENSAE Pierre Ndiaye}\\[2pt]
  {\small Impact des transferts de migrants sur la pauvreté multidimensionnelle des enfants au Sénégal}
\end{center}

\vspace{4pt}

\noindent\textbf{Principe.} Approche N-MODA (ANSD/UNICEF) : sept dimensions du
bien-être, déclinées en indicateurs binaires (privé / non privé). Un enfant est
\textbf{privé dans une dimension} s'il l'est dans \emph{au moins un} de ses
indicateurs (union intra-dimension). Il est \textbf{pauvre} (N-MODA) s'il est
privé dans \textbf{au moins $k=4$ dimensions sur 7}. Unité d'analyse~: l'enfant.
Trois groupes d'âge (\textbf{0-4}, \textbf{5-14}, \textbf{15-17} ans)
déterminent les indicateurs applicables. Statistiques descriptives pondérées
par le poids de sondage (\texttt{hhweight})~; effets causals pondérés par les
poids d'appariement du PSM.

\vspace{8pt}
\hrule
\vspace{8pt}

% ---------- 1. ASSAINISSEMENT ----------
\dimtitle{1. Assainissement}
\begin{tabularx}{\textwidth}{L{3.1cm} X L{3.3cm}}
\arrayrulecolor{entete}\toprule
\textbf{Indicateur} & \textbf{Définition (privé si\dots)} & \textbf{Variables (2018 / 2021)} \\
\midrule
Type de sanitaire non amélioré & le ménage utilise des toilettes non améliorées~: latrines SANPLAT, latrines dallées simples, fosse rudimentaire, toilettes publiques, aucune toilette, autre (codes 7 à 12). & \texttt{s11q55} / \texttt{s11q54} \\
\rowcolor{ligne}
Partage des toilettes & le ménage partage ses toilettes avec d'autres ménages (dénominateur~: ménages à installation privée). & \texttt{s11q56} / \texttt{s11q55} \\
\bottomrule
\end{tabularx}

% ---------- 2. EAU ----------
\dimtitle{2. Eau}
\begin{tabularx}{\textwidth}{L{3.1cm} X L{3.3cm}}
\arrayrulecolor{entete}\toprule
\textbf{Indicateur} & \textbf{Définition (privé si\dots)} & \textbf{Variables (2018 / 2021)} \\
\midrule
Source d'eau non améliorée & la source de boisson est non améliorée --- puits ouvert cour/concession (5), puits ouvert ailleurs (6), source non aménagée (12), fleuve/rivière/lac (13), vendeur ambulant (16), autre (17), à l'une ou l'autre saison --- \emph{et} le ménage ne traite pas son eau (filtre de traitement $\neq$ « oui »). & Source~: \texttt{s11q27a}, \texttt{s11q27b} / \texttt{s11q26a}, \texttt{s11q26b}. Filtre traitement~: \texttt{s11q32} / \texttt{s11q31} \\
\rowcolor{ligne}
Temps d'accès à l'eau & le temps de collecte (aller + attente à la source) dépasse 30 min à l'une ou l'autre saison. & \texttt{s11q29a}, \texttt{s11q31a} (+ heures/minutes) / \texttt{s11q28a}, \texttt{s11q30a} \\
\bottomrule
\end{tabularx}

% ---------- 3. LOGEMENT ----------
\dimtitle{3. Logement}
\begin{tabularx}{\textwidth}{L{3.1cm} X L{3.3cm}}
\arrayrulecolor{entete}\toprule
\textbf{Indicateur} & \textbf{Définition (privé si\dots)} & \textbf{Variables (2018 / 2021)} \\
\midrule
Débarras des ordures & le ménage évacue ses ordures de façon inadéquate~: brûlées (3), dépotoir sauvage (5), autre (6). & \texttt{s11q54} / \texttt{s11q53} \\
\rowcolor{ligne}
Surpeuplement & le ménage compte \textbf{au moins 4 personnes par pièce} (ratio taille/pièces $\geq 4$~; lecture ANSD de « plus de 3 personnes par pièce »). & \texttt{hhsize} (roster) $/$ \texttt{s11q02} (nombre de pièces) \\
\bottomrule
\end{tabularx}

% ---------- 4. NUTRITION ----------
\dimtitle{4. Nutrition}
\begin{tabularx}{\textwidth}{L{3.1cm} X L{3.3cm}}
\arrayrulecolor{entete}\toprule
\textbf{Indicateur} & \textbf{Définition (privé si\dots)} & \textbf{Variables (2018 / 2021)} \\
\midrule
Insécurité alimentaire (FIES) & un membre du ménage a, faute de ressources, sauté un repas, mangé moins que nécessaire, manqué de nourriture, eu faim sans manger, ou passé une journée sans manger. \emph{La diversité des repas n'est pas retenue (module non comparable entre vagues).} & \texttt{s08aq04} à \texttt{s08aq08} \\
\bottomrule
\end{tabularx}

% ---------- 5. SANTE ----------
\dimtitle{5. Santé \small\normalfont(dimension = union des deux indicateurs)}
\begin{tabularx}{\textwidth}{L{3.1cm} X L{3.3cm}}
\arrayrulecolor{entete}\toprule
\textbf{Indicateur} & \textbf{Définition (privé si\dots)} & \textbf{Variables (2018 / 2021)} \\
\midrule
Combustible solide pour cuisiner & le ménage cuisine avec un combustible solide~: bois (ramassé/acheté), charbon de bois, déchets animaux, autres. & \texttt{s11q53\_\_1,2,3,7,8} / \texttt{s11q52} \\
\rowcolor{ligne}
Accès à une structure de santé & l'enfant ne peut accéder \textbf{à pied} à une structure de santé publique (hôpital, service 5~; autre centre de santé public, service 6), c.-à-d. le mode habituel pour rejoindre la plus proche n'est pas la marche. & Module communautaire \texttt{s02\_co}~: \texttt{s02q00}, \texttt{s02q01\_\_5/\_\_6}, \texttt{s02q02} \\
\bottomrule
\end{tabularx}

% ---------- 6. PROTECTION ----------
\dimtitle{6. Protection de l'enfant}
\begin{tabularx}{\textwidth}{L{3.1cm} X L{3.3cm}}
\arrayrulecolor{entete}\toprule
\textbf{Indicateur} & \textbf{Définition (privé si\dots)} & \textbf{Variables (2018 / 2021)} \\
\midrule
Absence d'acte de naissance & l'enfant (\textbf{$<$ 15 ans}) ne dispose pas d'un acte de naissance. & \texttt{s01q05} \\
\rowcolor{ligne}
Travail des enfants & l'enfant (\textbf{5-14 ans}) effectue un travail économique \emph{ou} domestique d'au moins 1 heure. Domestique~: heures de courses, travaux domestiques, garde, corvées d'eau et de bois. Économique~: travail rémunéré ou pour propre compte. & Domestique~: \texttt{s04q01} à \texttt{s04q05}. Économique~: \texttt{s04q06} à \texttt{s04q09} \\
Séparation parentale & l'enfant ne vit pas avec ses \textbf{deux} parents biologiques (au moins un parent n'habite pas le ménage). & \texttt{s01q22} (père) et \texttt{s01q29} (mère) \\
\bottomrule
\end{tabularx}

% ---------- 7. EDUCATION ----------
\dimtitle{7. Éducation}
\begin{tabularx}{\textwidth}{L{3.1cm} X L{3.3cm}}
\arrayrulecolor{entete}\toprule
\textbf{Indicateur} & \textbf{Définition (privé si\dots)} & \textbf{Variables (2018 / 2021)} \\
\midrule
Non-scolarisation & l'enfant (\textbf{5-14 ans}) n'est pas scolarisé. & \texttt{scol} (base individus) \\
\rowcolor{ligne}
Illettrisme & l'enfant (\textbf{15-17 ans}) ne sait ni lire ni écrire. & \texttt{alfab} / \texttt{alfa} (base individus) \\
NEET & \emph{Indicateur écarté.} La valeur ANSD (85,7\,\%) correspond à la seule absence d'emploi, sans la condition de non-scolarisation, et compte donc les élèves comme privés~: inutilisable. & (\texttt{scol}, \texttt{activ7j}) \\
\bottomrule
\end{tabularx}

\vspace{10pt}
\hrule
\vspace{6pt}
{\small\itshape Note~: variables 2021 précisées après le « / » lorsqu'elles
diffèrent de 2018. Les âges indiqués en gras restreignent l'indicateur au
groupe concerné~; ailleurs l'indicateur, mesuré au niveau du ménage ou de la
localité, s'applique à tous les groupes d'âge.}

\end{document}
